\documentclass[11pt,a4paper]{article}

\usepackage[utf8]{inputenc}
\usepackage[T1]{fontenc}
\usepackage{amsmath}
\usepackage{amsfonts}
\usepackage{amssymb}
\usepackage{geometry}
\usepackage{hyperref}

\geometry{margin=1in}

\title{Theoretical Foundations of VarShare: Regularized Task-Specific Residuals}
\author{}
\date{}

\begin{document}

\maketitle

\section{Core Architecture}
The ``Deterministic VarShare'' ablation operates purely on the mean weights, removing variational sampling. The formulation is a shared backbone with additive, regularized, task-specific mean shifts (residuals):
\[ w_t = \theta + \mu_t \]
Where:
\begin{itemize}
    \item $\theta$: Shared parameters optimized across all tasks.
    \item $\mu_t$: Task-specific residuals optimizing only for task $t$.
    \item \textbf{Regularization}: An L2 penalty (or KL divergence equivalent) is applied to $\mu_t$, explicitly penalizing its magnitude and forcing it towards zero unless absolutely necessary.
\end{itemize}

\section{Theoretical Advantages}

\subsection{Accelerated Optimization via Constructive Interference}
\begin{itemize}
    \item \textbf{Mechanism}: The shared backbone $\theta$ receives gradient updates from \textit{all} tasks simultaneously. 
    \item \textbf{Benefit}: If there are $N$ tasks, early layers in $\theta$ converge on general, universal features (e.g., edge detection, basic state-space correlations) $N$ times faster than independent networks, drastically improving sample efficiency.
\end{itemize}

\subsection{Mitigation of Negative Transfer (The ``Pressure Valve'')}
\begin{itemize}
    \item \textbf{Mechanism}: In fully shared networks, conflicting task gradients cancel each other out, leading to oscillation and ``Negative Transfer'' where the network performs poorly on all tasks.
    \item \textbf{Benefit}: The regularized residuals $\mu_t$ act as structural shock absorbers. If tasks agree, gradients sum into $\theta$. If tasks disagree violently, the gradients for $\theta$ sum to zero (stopping oscillation), and the gradients flow directly into the task-specific $\mu_t$'s. This allows the network to automatically route ``agreements'' into the shared backbone and ``disagreements'' into the residuals.
\end{itemize}

\subsection{Information Bottlenecks and Improved Generalization}
\begin{itemize}
    \item \textbf{Mechanism}: The heavy L2 penalty on $\mu_t$ makes task-specific parameter growth ``expensive.''
    \item \textbf{Benefit}: Tasks are strongly incentivized to reuse the shared logic in $\theta$ rather than memorizing noise independently. This acts as an implicit regularizer against overfitting on individual tasks, promoting generalization.
\end{itemize}

\subsection{Parameter Efficiency}
\begin{itemize}
    \item \textbf{Mechanism}: Instead of holding $N$ full models in memory.
    \item \textbf{Benefit}: The architecture scales at roughly $\mathcal{O}(P + N \times p_{residual})$. Using low-rank variants (VarShare-LoRA), 50 tasks can be trained with only a marginal increase in parameter count compared to a single baseline model.
\end{itemize}

\section{Disadvantages and Failure Modes}

\subsection{The ``Entanglement'' Problem}
\begin{itemize}
    \item \textbf{Issue}: Because every task still mathematically passes through exactly the same shared $\theta$, $\mu_t$ only provides an additive shift. If two tasks require fundamentally contradictory feature transformations, $\theta$ gets caught in the middle. The network is forced to pay a massive L2 penalty to let the $\mu$ values grow large enough to completely ``undo'' $\theta$.
    \item \textbf{Contrast}: Explicit routing methods (like Soft Modularization) completely isolate parameters, avoiding this entanglement entirely.
\end{itemize}

\subsection{Hyper-Sensitivity to Regularization (Tuning Bottleneck)}
\begin{itemize}
    \item \textbf{Issue}: The entire architecture relies on the exact balance of the L2 penalty (the KL Beta parameter). 
    \begin{itemize}
        \item If too high: $\mu \to 0$. The architecture collapses into a rigid Shared Bottom network (causing massive negative transfer).
        \item If too low: $\theta$ growth stalls. Tasks find it cheaper to just use $\mu_t$ independently. The architecture collapses into $N$ Independent Networks (losing all sample-efficiency benefits).
    \end{itemize}
\end{itemize}

\subsection{Forced Late-Stage Sharing (The ``Late Divergence'' Problem)}
\begin{itemize}
    \item \textbf{Issue}: Standard application of VarShare forces the network to maintain the shared $\theta$ bottleneck even at the deepest layers of the network. If tasks share low-level features but have completely divergent goals at the high level, forcing them through a shared deep layer is computationally wasteful and mathematically unnatural.
\end{itemize}

\subsection{Limited Geometric Expressivity (Linear Shifts)}
\begin{itemize}
    \item \textbf{Issue}: $w_t = \theta + \mu_t$ is a purely linear, additive shift in weight space. It struggles to efficiently represent complex, non-linear geometric transformations (like rotating a feature space) required to fix a shared representation that conflicts with a specific task.
\end{itemize}

\section{Amplifying the Strengths (Advanced Strategies)}
To maximize the performance of the regularized-residual architecture without introducing unfair advantages or unprincipled heuristics, we can deploy the following advanced strategies:

\subsection{Adversarial Representation Alignment (ARA)}
\begin{itemize}
    \item \textbf{Goal}: Maximize the transferability and general usefulness of the shared backbone $\theta$.
    \item \textbf{Concept}: Attach an ``Adversary'' network to the output of $\theta$ that tries to predict the active task ID. Train $\theta$ via a Gradient Reversal Layer to explicitly fool the Adversary.
    \item \textbf{Why it works}: It mathematically forbids $\theta$ from learning lazy, task-specific shortcuts (e.g., memorizing colors to identify a task). It forces $\theta$ to extract only universal, foundational physics/logic, perfectly delegating $100\%$ of the task-specific routing burden precisely onto the task residuals $\mu_t$.
\end{itemize}

\subsection{L1 Regularization (Sparse Residuals)}
\begin{itemize}
    \item \textbf{Goal}: Supercharge the ``Shock Absorber'' effect without needlessly shifting weights that already agree.
    \item \textbf{Concept}: Swap (or combine) the standard L2 penalty on $\mu_t$ for an L1 penalty, which heavily induces sparsity.
    \item \textbf{Why it works}: While L2 shrinks all numbers, L1 forces the vast majority of the $\mu_t$ matrix to be exactly $0.0$. This ensures the task strictly relies on $\theta$ for almost everything, only allowing surgical alterations to the tiny fraction of weights causing catastrophic task conflicts.
\end{itemize}

\subsection{VarShare-LoRA (Low-Rank Residuals)}
\begin{itemize}
    \item \textbf{Goal}: Physically enforce the Information Bottleneck and drastically improve parameter efficiency.
    \item \textbf{Concept}: Constrain the residual to a low-rank factorization: $\mu_t = A \times B$ where the inner rank is very small.
    \item \textbf{Why it works}: It structurally prevents the task-specific head from having enough parameter capacity to memorize the task independently. The network has no choice but to funnel its learning capacity into the shared backbone $\theta$, maximizing constructive interference.
\end{itemize}

\section{Feasible Solutions and Future Architectural Ablations}
To address the fundamental weaknesses of the regularized-residual approach, the following architectural ablations present the most realistic and computationally feasible paths forward for future research:

\subsection{Solving Entanglement: Soft Muting / Selective Gating}
\begin{itemize}
    \item \textbf{Addresses}: The Entanglement Problem (3A)
    \item \textbf{Concept}: Introduce a task-specific vector of learned gates $g_t$ (passed through a Sigmoid so values are strictly between 0 and 1). The forward pass becomes: $y_{task} = (\text{sigmoid}(g_t) \odot \theta + \mu_t)x$.
    \item \textbf{Why it works}: If a task fundamentally disagrees with a specific feature in $\theta$, it learns $g_{t,i} = 0.0$. By multiplying the conflicting feature by zero, it entirely silences it \textit{without} having to learn a massive negative number in $\mu_t$, perfectly avoiding the L2 penalty and disentangling the task.
\end{itemize}

\subsection{Solving Hyper-Sensitivity to Regularization}
\begin{itemize}
    \item \textbf{Addresses}: Hyper-Sensitivity / Tuning Bottleneck (3B)
    \item \textbf{Solution B1: L1 Regularization (Sparse Residuals)}
    \begin{itemize}
        \item Swap the standard L2 penalty on $\mu_t$ for an L1 penalty (or ElasticNet). L1 forces the vast majority (e.g., $95\%$) of the $\mu_t$ matrix to be exactly $0.0$. This ensures the task strictly relies on $\theta$ for almost everything, only acting as a ``shock absorber'' for the exact $5\%$ of weights causing catastrophic task conflicts.
    \end{itemize}
    \item \textbf{Solution B2: PCGrad (Projecting Conflicting Gradients)}
    \begin{itemize}
        \item Wrap the optimizer with PCGrad. If Task A and Task B's gradients point in opposite directions (cosine similarity $< 0$) regarding $\theta$, PCGrad mathematically projects them to be orthogonal. This manually removes interference, meaning $\theta$ never gets caught in a tug-of-war, allowing the L2 penalty on $\mu_t$ to safely stay high without risking negative transfer.
    \end{itemize}
    \item \textbf{Solution B3: Learned Hyperprior}
    \begin{itemize}
        \item Make the regularization coefficient $\beta_t$ a learned parameter per task, parameterized as $\text{Softplus}(\beta_t)$. While tricky to balance in RL, the network dynamically finds its own optimal sharing ratio per task without manual hyperparameter search.
    \end{itemize}
\end{itemize}

\subsection{Solving Late-Stage Sharing: Progressive Depth Decay}
\begin{itemize}
    \item \textbf{Addresses}: Forced Late-Stage Sharing (3C)
    \item \textbf{Concept}: Apply a fixed hyperparameter schedule across the network's depth. The L2 penalty for a layer is multiplied by a decay factor (either linear or exponential) based on depth. (e.g., Layer 1: 1.0x penalty, Layer 2: 0.5x penalty, Final Layer: 0.0x penalty).
    \item \textbf{Why it works}: It natively forces the network to share raw feature extraction at the bottom (where a ``middle ground'' is mathematically optimal), while allowing the final output heads to be completely independent without paying any regularization penalty (as final action decoders rarely share cohesive task logic).
\end{itemize}

\subsection{Solving Geometric Expressivity: VarShare-FiLM}
\begin{itemize}
    \item \textbf{Addresses}: Limited Geometric Expressivity (3D)
    \item \textbf{Concept}: Instead of modifying weights additively, allow the task to modify the activations multiplicatively. After $\theta$ processes the input $x$, apply a task-specific learned scale vector $\gamma_t$ and shift vector $\beta_t$: $y_{task} = \gamma_t \odot (\theta x) + \beta_t$.
    \item \textbf{Why it works}: Addition is mathematically inefficient for transforming geometry (like rotating a feature space). Multiplication allows the task to dynamically stretch, compress, or invert the shared feature space. It is orders of magnitude more parameter-efficient than an additive $\mu_t$ matrix and vastly more expressive.
\end{itemize}

\end{document}
